\documentclass[11pt]{amsart}
\usepackage{amsmath,amssymb,amsthm}
\usepackage[margin=1.2in]{geometry}
\usepackage{newtxtext,newtxmath}
\usepackage{graphicx}
\usepackage{microtype}
\usepackage[colorlinks=true, linkcolor=blue, citecolor=blue, urlcolor=blue]{hyperref}

\newtheorem{lemma}{Lemma}[section]
\newtheorem{remark}[lemma]{Remark}
\newtheorem*{theoremA}{Theorem}

\newcommand{\Z}{\mathbb{Z}}
\newcommand{\llog}{\log\log}
\newcommand{\lllog}{\log\log\log}

\begin{document}

\title[Constructing Carmichael numbers deterministically]{A deterministic quasi-polylogarithmic algorithm\\ for constructing Carmichael numbers}

\author{Dan Brown$^*$}
\email{jdb1729@gmail.com}
\thanks{$^*$Author of \emph{The Da Vinci Code}.}
\thanks{In memory of Abraham Anthony ``Mickey'' Br\"umleve (August 24, 1986 -- December 5, 2019): A thinker, a visionary, and a source of fantastic ideas.}

\date{August 24, 2026}

\begin{abstract}
We describe a deterministic algorithm that, given a large integer $n$,
outputs a Carmichael number $m \in (n, n^{1+o(1)}]$ together with its
complete prime factorization, in time
$\exp\bigl(O(\llog n \cdot \lllog n)\bigr) = (\log n)^{O(\lllog n)}$.
The correctness proof uses no unproved hypothesis and no randomness. By
contrast, the fastest known deterministic algorithm producing a proven
prime greater than $n$ takes time $n^{1/2+o(1)}$. The construction
assembles effective results of Alford, Granville and Pomerance with a
subset-product search carried out by dynamic programming over the group
$(\Z/L\Z)^\times$ for a highly composite modulus $L$ of quasi-polylogarithmic
size. The output certificate is verifiable in polylogarithmic time.
\end{abstract}

\maketitle

\section{Introduction}

Given $n$, how quickly can one produce, with proof, a number larger than
$n$ of a prescribed multiplicative type? For primes this is a well-known
open problem: a random search succeeds in expected polylogarithmic time,
but the best deterministic bound is $n^{1/2+o(1)}$, obtained by binary
search against the prime counting function computed by the analytic
method of Lagarias and Odlyzko \cite{LO}; see the Polymath project
\cite{Pol} for the state of the art and for the obstructions near the
exponent $1/2$.

This note concerns Carmichael numbers: composite $m$ such that
$a^m \equiv a \pmod{m}$ for every integer $a$. By Korselt's criterion
\cite{Kor}, $m$ is Carmichael if and only if $m$ is composite, squarefree,
and $p-1 \mid m-1$ for every prime $p \mid m$. The criterion refers only
to the prime factors of $m$, and those factors can be small. This opens a
route that is unavailable for primes: build $m$ as a product of many
small primes chosen so that Korselt's criterion holds automatically. The
existence machinery is exactly that of Alford, Granville and Pomerance
\cite{AGP} in their proof that there are infinitely many Carmichael
numbers. Our observation is that, for the parameters relevant here, every
object in their proof is small enough to be found by exhaustive
deterministic search, and their existence statements are effective.

\begin{theoremA}
There is a deterministic algorithm and a constant $n_0$, computable in
principle, such that for every $n \ge n_0$ the algorithm outputs an
integer $m \in (n, n^{1+o(1)}]$ together with its prime factorization,
$m$ is a Carmichael number, and the running time is
\[
\exp\bigl(O(\llog n \cdot \lllog n)\bigr).
\]
The factorization constitutes a certificate of the Carmichael property
verifiable in deterministic polylogarithmic time.
\end{theoremA}

To read the time bound: the input has $\log n$ bits, so a
polylogarithmic-time algorithm (the ideal) would run in time
$(\log n)^{O(1)}$. Our bound exceeds that by only a factor of $\lllog n$
in the exponent, and it is smaller than $\exp((\log n)^{c})$ for every $c>0$ ---
in particular smaller than any fixed positive power of $n$. We are not
aware of any previously published deterministic time bound, at
any subexponential level, for producing a Carmichael number beyond a
given bound.

An outline of the construction is as follows. Using
Lemma~\ref{lem:smooth} we build a modulus $L$, a product of a few dozen
small primes $q$ chosen so that every $q-1$ has only very small prime
factors. Using Lemma~\ref{lem:pigeon} we find a multiplier $k$ such that
many numbers of the form $dk+1$, with $d$ running over divisors of $L$,
are prime; such primes $p$ satisfy $p - 1 \mid kL$ by design. Any
product $m$ of distinct such primes with $m \equiv 1 \pmod{kL}$ is then
automatically a Carmichael number by Korselt's criterion, and
Lemma~\ref{lem:vebk} guarantees that products with $m \equiv 1$ are
plentiful enough for an exhaustive search to find them until their
accumulated product exceeds $n$.

Related work: the recognition problem (deciding whether a
given $n$ is Carmichael) was recently shown to be in P \cite{SW};
large-scale computations of Carmichael numbers by subset products
\cite{LN, AGHS} use the same skeleton as our algorithm with randomized
and heuristic searches in place of the exhaustive steps; and Larsen
\cite{Lar} proved an effective analogue of Bertrand's postulate for
Carmichael numbers, discussed in Remark~\ref{rem:larsen}.

Throughout, $\ell_2 = \llog n$ and $\ell_3 = \lllog n$ (natural
logarithms), $P^+(k)$ is the largest prime factor of $k$, and running
time counts bit operations. $C_1$ denotes an effective absolute constant.

\section{Components}

All four inputs below are unconditional, and effective in the sense that
their constants are computable in principle.

\begin{lemma}[Smooth shifted primes; \cite{AGP}, Theorem 3 and p.~707--708]
\label{lem:smooth}
Let $\pi(x,y) = \#\{p \le x \text{ prime} : P^+(p-1) \le y\}$. For every
$E < 5/12$ there are computable constants $\gamma(E) > 0$ and $x_1(E)$
such that $\pi(x, x^{1-E}) \ge \gamma(E)\,\pi(x)$ for all $x \ge x_1(E)$.
\end{lemma}

That is, a definite fraction of the primes $p$ up to $x$ have the
property that $p-1$ factors entirely into primes of size at most
$x^{2/3}$ (taking $E = 1/3$, which we fix from now on; write
$\gamma = \gamma(1/3)$). We use such primes as the building blocks of
our modulus $L$ because the smoothness of the $q-1$ keeps the group
exponent $\lambda(L)$ small, which is what makes
Lemma~\ref{lem:vebk} below effective with very short sequences.

\begin{lemma}[Pigeonhole over shifts; \cite{AGP}, Theorem 3.1]
\label{lem:pigeon}
For each $B$ in the class $\mathcal{B}$ of \cite{AGP} there are computable
constants $z_3(B)$ and $D_B$ with the following property. If
$x > z_3(B)$, if $L$ is squarefree with no prime factor exceeding
$x^{(1-B)/2}$, and if $\sum_{q \mid L} 1/q \le (1-B)/32$ (sum over
primes), then there is a positive integer $k \le x^{1-B}$ with
$(k,L)=1$ such that
\[
\#\{d \mid L : dk+1 \le x,\ dk+1 \text{ prime}\}
\;\ge\; \frac{2^{-D_B-2}}{\log x}\,\#\{d \mid L : d \le x^{B}\}.
\]
Moreover every $B < 5/12$ lies in $\mathcal{B}$, effectively.
\end{lemma}

We fix $B = 2/5$ and write $D = D_{2/5}$. Numbers of the form $dk+1$
with $d \mid L$ have no reason to avoid being prime, and the lemma
confirms that for at least one small multiplier $k$ a decent proportion
of them (one in roughly $\log x$, up to a constant) are indeed prime. These are the primes our Carmichael number will be made
of: each satisfies $p - 1 = dk$, which divides $kL$, and that is
exactly the divisibility Korselt's criterion requires.

\begin{lemma}[van Emde Boas--Kruyswijk \cite{EK}; see \cite{AGP}, Theorem 1.1]
\label{lem:vebk}
Let $G$ be a finite abelian group whose maximal element order is $m^*$.
Any sequence of more than $m^*\bigl(1 + \log(|G|/m^*)\bigr)$ elements of
$G$ contains a nonempty subsequence whose product is the identity.
\end{lemma}

What is striking here is how few elements suffice for a sub-collection
with product equal to the identity: not on the order of the group size
$|G|$, but on the order of the largest order $m^*$ of a single element,
times a logarithm. For our group $G = (\Z/L\Z)^\times$ the size $|G|$ is huge
but $m^*$ divides $\lambda(L)$, which the smoothness in
Lemma~\ref{lem:smooth} keeps tiny; this gap is the engine of the whole
construction, exactly as in \cite{AGP}.

\begin{lemma}[AKS \cite{AKS}]
\label{lem:aks}
Primality of an $N$-bit integer is decidable deterministically in time
polynomial in $N$.
\end{lemma}

\section{The algorithm}
\label{sec:algo}

Input: $n \ge n_0$.

\begin{enumerate}
\item[\textbf{1.}] \textbf{Scales.} Set
$z = \lceil C_1 \ell_2 \ell_3 \rceil$ with
$C_1 = \max(16/\gamma,\,10^3)$, then $y = \lceil z^{2/3} \rceil$ and
$T = \lceil 3\ell_2 \rceil$.

\item[\textbf{2.}] \textbf{Moduli.} By a sieve up to $z$ and trial
division of each $q-1$, list the primes
$q \in (\sqrt{z},\, z]$ with $P^+(q-1) \le y$. Let $Q$ be the $T$
largest, $L = \prod_{q \in Q} q$, and $x = \lceil L^5 \rceil$.

\item[\textbf{3.}] \textbf{Shift scan.} For $k = 1, 2, \dots$ with
$(k, L) = 1$: enumerate the $2^T$ divisors $d$ of $L$ as subset products
of $Q$, and test each $dk+1 \le x$ for primality
(Lemma~\ref{lem:aks}). Let
$P_k = \{dk+1 : d \mid L,\ dk+1 \le x,\ dk+1 \text{ prime},\ dk+1 > z\}$.
Stop at the first $k$ with $|P_k| \ge (\log n)^{1.2}$; write
$P = P_k$.

\item[\textbf{4.}] \textbf{Extraction.} Compute
$\lambda(L) = \operatorname{lcm}\{q - 1 : q \in Q\}$ (each $q-1$ is
fully factored) and set $N^* = \lceil \lambda(L)(1 + \log L)\rceil + 1$.
Set $m = 1$. While $m \le n$: take the next $N^*$ unused elements of
$P$; by dynamic programming over residues modulo $L$ (process the
elements one at a time, keep one witness subset per reachable residue,
and use a per-element snapshot so that each element is used at most
once) find a nonempty subset $S$ with
$\prod_{p \in S} p \equiv 1 \pmod{L}$; replace $m$ by
$m \prod_{p \in S} p$ and mark the elements of $S$ used.

\item[\textbf{5.}] \textbf{Output and certificate.} Output $m$ and the
list of primes used. Verify: the factors are distinct and prime, their
product is $m$, the count of factors is at least $3$, and
$p - 1 \mid m - 1$ for each factor $p$.
\end{enumerate}

Three of the five steps are searches, each guaranteed a target ---
Lemma~\ref{lem:smooth} for step 2, Lemma~\ref{lem:pigeon} for step 3,
Lemma~\ref{lem:vebk} for step 4 --- so for $n \ge n_0$ no link of the
kill chain comes up empty.

\section{Proof of the theorem}

\begin{lemma}[Step 2 succeeds]
For $n$ large, at least $T$ primes $q \in (\sqrt z, z]$ satisfy
$P^+(q-1) \le y$.
\end{lemma}

\begin{proof}
Apply Lemma~\ref{lem:smooth} at scale $z$ (this requires
$z \ge x_1(1/3)$, absorbed into $n_0$): at least $\gamma\,\pi(z)$ primes
$q \le z$ have $P^+(q-1) \le z^{2/3} \le y$, so at least
$\gamma\,\pi(z) - \pi(\sqrt z) \ge \tfrac{\gamma}{2}\pi(z)$ of them lie
in $(\sqrt z, z]$. Since $\pi(z) \ge z/\log z \ge (C_1/2)\,\ell_2$ for
large $n$, this count is at least
$\tfrac{\gamma}{2}\cdot\tfrac{C_1}{2}\,\ell_2 = \tfrac{\gamma C_1}{4}\,\ell_2
\ge 4\ell_2 \ge T$, using $C_1 \ge 16/\gamma$.
\end{proof}

\begin{lemma}[Step 3 halts]
\label{lem:halt}
The scan accepts some $k \le x^{3/5}$.
\end{lemma}

\begin{proof}
We check the hypotheses of Lemma~\ref{lem:pigeon} for $B = 2/5$. $L$ is
squarefree and its prime factors are at most $z$, while
$x^{(1-B)/2} = x^{0.3} \ge 2^{1.5\,T} > z$ (as $L \ge 2^T$). Also
$\sum_{q \in Q} 1/q \le T/\sqrt z \to 0$, so the sum is at most
$(1-B)/32 = 0.01875$ for large $n$; and $x > z_3(2/5)$ for large $n$.
Lemma~\ref{lem:pigeon} yields $k \le x^{3/5}$, $(k,L) = 1$, with
\[
\#\{d \mid L : dk+1 \le x \text{ prime}\}
\ \ge\ \frac{2^{-D-2}}{\log x}\; \#\{d \mid L : d \le x^{2/5}\}.
\]
Every divisor of $L$ is at most $L \le x^{2/5}$, so the last factor
equals $2^{T} \ge (\log n)^{3\log 2} = (\log n)^{2.07\ldots}$. Since
$2^{D+2}\log x = O(\ell_2 \ell_3)$, the right side is
$(\log n)^{2.07 - o(1)}$; discarding the at most $\pi(z)$ elements with
$dk+1 \le z$ leaves more than $(\log n)^{1.2}$ for large $n$. The scan
tests precisely this count, so it accepts at this $k$ if it has not
stopped earlier.
\end{proof}

\begin{lemma}[Step 4 always finds a subset, and the pool suffices]
\label{lem:pool}
Each round of step 4 finds a nonempty $S$ with
$\prod_{p\in S} p \equiv 1 \pmod L$, and the loop ends before $P$ is
exhausted.
\end{lemma}

\begin{proof}
The elements of $P$ are distinct primes (distinct $d$ give distinct
$dk+1$) exceeding $z \ge \max Q$, hence coprime to $L$; they are
elements of $G = (\Z/L\Z)^\times$. The maximal order of an element of
$G$ divides $\lambda(L)$, and
$N^* > \lambda(L)(1 + \log(|G|/\lambda(L)))$, so by
Lemma~\ref{lem:vebk} any $N^*$ elements contain a nonempty subsequence
with product $1$. The dynamic program records a witness subset for every
residue reachable by a nonempty subset of the current pool, so it
recovers such an $S$.

For the pool size: each $q - 1$ ($q \in Q$) is $y$-smooth, so
$\lambda(L)$ divides $\prod_{p \le y} p^{\lfloor \log_p z\rfloor}$ and
$\log \lambda(L) \le \pi(y)\log z = O(z^{2/3}) = O((\ell_2\ell_3)^{2/3})$,
whence $N^* = \exp\bigl(O((\ell_2\ell_3)^{2/3})\bigr) = (\log n)^{o(1)}$.
Each round multiplies $m$ by a product congruent to $1$ modulo $L$ and
larger than $1$, hence at least $L + 1$; since each $q > \sqrt z$ gives
$\log L \ge \tfrac{T}{2}\log z \ge \ell_2 \ell_3$ for large $n$, the
loop ends after at most $\lceil \log n / \log(L+1) \rceil \le
\log n/(\ell_2\ell_3)$ rounds. Total consumption is at most
$N^* \log n/(\ell_2 \ell_3) = (\log n)^{1+o(1)}/(\ell_2\ell_3)
< (\log n)^{1.2} \le |P|$.
\end{proof}

\begin{lemma}
The output $m$ is a Carmichael number, and $m \in (n, n^{1+o(1)}]$.
\end{lemma}

\begin{proof}
$m > n$ by the loop condition. The subsets $S_i$ are disjoint and their
elements are distinct primes, so $m$ is squarefree. Every factor $p$
satisfies $p = dk+1$ with $d \mid L$, so $p - 1 = dk \mid kL$. Each
$p \equiv 1 \pmod k$ gives $m \equiv 1 \pmod k$; each round's product is
$\equiv 1 \pmod L$, so $m \equiv 1 \pmod L$; and $(k, L) = 1$, so
$m \equiv 1 \pmod{kL}$. Hence $p - 1 \mid kL \mid m - 1$ for every
$p \mid m$. All factors are at most $x = e^{O(\ell_2\ell_3)}$, so $m$
has at least $\log n/\log x \ge 3$ factors for large $n$; in particular
$m$ is composite, and Korselt's criterion applies. Finally, before its
last round the loop had $m \le n$, and one round multiplies $m$ by at
most $x^{N^*} = \exp\bigl(O(\ell_2\ell_3)\,(\log n)^{o(1)}\bigr)
= n^{o(1)}$.
\end{proof}

\begin{lemma}[Running time]
The algorithm runs in time $\exp\bigl(O(\ell_2\ell_3)\bigr)$, and the
certificate of step 5 is verifiable in time $(\log n)^{2+o(1)}$.
\end{lemma}

\begin{proof}
Step 2 costs $\operatorname{poly}(z)$. Step 3 examines at most
$x^{3/5} = e^{3\log L}$ shifts, each costing
$2^T \operatorname{poly}(\log x)$ by Lemma~\ref{lem:aks}; in total
$e^{O(\ell_2\ell_3)}$. Each round of step 4 costs at most
$L \cdot N^*$ dictionary operations on $O(\log L)$-bit keys, i.e.
$e^{O(\ell_2\ell_3)}$, and there are at most $\log n$ rounds. Assembling
$m$ and checking the certificate involves $(\log n)^{1+o(1)}$ factors of
$(\log n)^{O(\ell_3)}$ bits against an $O(\log n)$-bit $m$, which is
$(\log n)^{2+o(1)}$ bit operations.
\end{proof}

This proves the theorem. \qed

\section{Remarks}

\begin{remark}[The threshold]
$n_0$ collects $x_1(E)$, $z_3(B)$, $D_B$ and the finitely many
``for $n$ large'' comparisons above. The constants of \cite{AGP} are
computable in principle; we are not aware of numerical values in the
literature, either for their Theorem 5 or for the constants used here.
Making $n_0$ numeric would require explicit
versions of the zero-density inputs to \cite{AGP}, along the lines of
the explicit Bombieri--Vinogradov literature.
\end{remark}

\begin{remark}[Short intervals]
\label{rem:larsen}
The algorithm overshoots $n$ by a factor $n^{o(1)}$ and cannot aim into
a window such as $(n, 2n]$. Larsen \cite{Lar} proved, effectively, that
$[x,\, x + x/(\log x)^{1/(2+\delta)}]$ contains Carmichael numbers for
large $x$, with building blocks of size
$\exp((\llog x)^{2+o(1)})$; his proof appears amenable to the same
treatment and would give a deterministic
$\exp((\llog n)^{2+o(1)})$-time construction inside such windows. We
have not carried this out.
\end{remark}

\begin{remark}[Comparison with primes]
Producing a proven prime exceeding $n$ in deterministic
polylogarithmic time is open even under the Riemann Hypothesis, and the
record $n^{1/2+o(1)}$ \cite{LO, Pol} has stood since 1987. The theorem
separates the two construction problems: the Carmichael property is a
conjunction of conditions on small prime factors, verifiable and
assemblable at scale $\exp(O(\ell_2 \ell_3))$, whereas primality of $m$
is a single condition at scale $m$.
\end{remark}

\begin{remark}[Practice]
A simplified variant (reservoir $\{q : q - 1 \mid \operatorname{lcm}(1..K)\}$,
no shift scan; the abundance of that reservoir is conjectural, which is
why the proof above uses the shifted reservoir of \cite{AGP}) produces
certified Carmichael numbers above $10^{150}$ in a few seconds of
interpreted code, e.g.\ a 219-digit example with 70 prime factors all
below $2 \cdot 10^{5}$. For $n = 10^{20}$ it returns the 65-digit
Carmichael number
\[
m = 32053309583031266791428889135359834328765733228091969901544344081
\]
with the 28 prime factors $13 \cdot 19 \cdot 23 \cdot 29 \cdot 31 \cdot
37 \cdot 43 \cdot 61 \cdot 67 \cdot 71 \cdot 73 \cdot 89 \cdot 127
\cdot 181 \cdot 199 \cdot 211 \cdot 281 \cdot 331 \cdot 397 \cdot 463
\cdot 617 \cdot 631 \cdot 991 \cdot 1321 \cdot 2311 \cdot 2521 \cdot
4621 \cdot 9241$. An implementation is available from the author.
\end{remark}

\begin{thebibliography}{9}

\bibitem{AGP} W.~R. Alford, A. Granville, C. Pomerance,
\emph{There are infinitely many Carmichael numbers},
Ann. of Math. (2) \textbf{140} (1994), 703--722.
\href{https://doi.org/10.2307/2118576}{doi:10.2307/2118576}.

\bibitem{AKS} M. Agrawal, N. Kayal, N. Saxena,
\emph{PRIMES is in P},
Ann. of Math. (2) \textbf{160} (2004), 781--793.
\href{https://doi.org/10.4007/annals.2004.160.781}{doi:10.4007/annals.2004.160.781}.

\bibitem{AGHS} W.~R. Alford, J. Grantham, S. Hayman, A. Shallue,
\emph{Constructing Carmichael numbers through improved subset-product
algorithms}, Math. Comp. \textbf{83} (2014), 899--915.
\href{https://doi.org/10.1090/S0025-5718-2013-02737-8}{doi:10.1090/S0025-5718-2013-02737-8}.

\bibitem{Kor} A. Korselt,
\emph{Probl\`eme chinois},
L'interm\'ediaire des math\'ematiciens \textbf{6} (1899), 142--143.

\bibitem{LO} J.~C. Lagarias, A.~M. Odlyzko,
\emph{Computing $\pi(x)$: an analytic method},
J. Algorithms \textbf{8} (1987), 173--191.
\href{https://doi.org/10.1016/0196-6774(87)90037-X}{doi:10.1016/0196-6774(87)90037-X}.

\bibitem{Lar} D. Larsen,
\emph{Bertrand's postulate for Carmichael numbers},
Int. Math. Res. Not. IMRN \textbf{2023}, no.~15, 13072--13098.
\href{https://doi.org/10.1093/imrn/rnac203}{doi:10.1093/imrn/rnac203}.

\bibitem{LN} G. L\"oh, W. Niebuhr,
\emph{A new algorithm for constructing large Carmichael numbers},
Math. Comp. \textbf{65} (1996), 823--836.
\href{https://doi.org/10.1090/S0025-5718-96-00692-8}{doi:10.1090/S0025-5718-96-00692-8}.

\bibitem{Pol} D.~H.~J. Polymath,
\emph{Deterministic methods to find primes},
Math. Comp. \textbf{81} (2012), 1233--1246.
\href{https://doi.org/10.1090/S0025-5718-2011-02542-1}{doi:10.1090/S0025-5718-2011-02542-1}.

\bibitem{SW} A. Shallue, J. Webster,
\emph{Algorithms for Carmichael numbers},
\href{https://arxiv.org/abs/2506.09903}{arXiv:2506.09903} (2025).

\bibitem{EK} P. van Emde Boas, D. Kruyswijk,
\emph{A combinatorial problem on finite abelian groups III},
Report ZW-1969-008, Mathematisch Centrum, Amsterdam, 1969.
\href{https://ir.cwi.nl/pub/7217}{ir.cwi.nl/pub/7217}.

\end{thebibliography}
\end{document}
